% RawHDR Composer 使用説明書
% 組版: lualatex manual.tex（2 回）。画像は fig/ に置く（capture.sh で撮った画面と、Camera Raw で現像した作例）。
\documentclass[paper=a4paper,fontsize=11pt,report,jafontsize=11pt]{jlreq}
\usepackage[hiragino-pron]{luatexja-preset}
\usepackage{graphicx}
\usepackage{tikz}
\usetikzlibrary{arrows.meta,calc,positioning,shapes.callouts}
\usepackage[most]{tcolorbox}
\usepackage{booktabs,tabularx,array,multirow}
\usepackage{fontawesome5}
\usepackage{xcolor}
\usepackage{enumitem}
\usepackage[hidelinks,pdfusetitle]{hyperref}
\usepackage{bookmark}

\hypersetup{pdftitle={RawHDR Composer 使用説明書},pdfauthor={RawHDR Composer}}

% ---- 色と囲み -----------------------------------------------------------------------
\definecolor{accent}{HTML}{C62828}
\definecolor{navy}{HTML}{1F3A5F}
\definecolor{tipbg}{HTML}{E8F4EA}
\definecolor{tipfr}{HTML}{2E7D32}
\definecolor{warnbg}{HTML}{FFF4E5}
\definecolor{warnfr}{HTML}{E65100}
\definecolor{notebg}{HTML}{EAF1FB}
\definecolor{notefr}{HTML}{1565C0}

\newtcolorbox{tip}[1][ヒント]{colback=tipbg,colframe=tipfr,fonttitle=\bfseries\sffamily,title={\faLightbulb[regular]\ #1},boxrule=0.6pt,arc=2pt,left=6pt,right=6pt}
\newtcolorbox{warn}[1][注意]{colback=warnbg,colframe=warnfr,fonttitle=\bfseries\sffamily,title={\faExclamationTriangle\ #1},boxrule=0.6pt,arc=2pt,left=6pt,right=6pt}
\newtcolorbox{memo}[1][メモ]{colback=notebg,colframe=notefr,fonttitle=\bfseries\sffamily,title={\faInfoCircle\ #1},boxrule=0.6pt,arc=2pt,left=6pt,right=6pt}
% 手順の見出し
\newtcolorbox{step}[2]{breakable,colback=white,colframe=navy,fonttitle=\bfseries\sffamily\large,title={ステップ #1\quad #2},boxrule=0.8pt,arc=2pt,left=6pt,right=6pt}

% ---- 画面の画像に注釈を描く ---------------------------------------------------------
% \shot[幅]{画像}{注釈}: 注釈の座標は (横, {1-縦})。横・縦とも画像の左上を 0、右下を 1 とする割合。
\tikzset{
  num/.style={circle,fill=accent,text=white,font=\bfseries\sffamily\small,inner sep=0pt,minimum size=1.35em,draw=white,line width=0.8pt},
  hl/.style={draw=accent,line width=1.6pt,rounded corners=2pt},
  arr/.style={-{Stealth[length=3.2mm,width=2.6mm]},line width=1.6pt,accent},
  lbl/.style={fill=white,draw=accent,line width=0.8pt,rounded corners=2pt,font=\sffamily\footnotesize,align=left,inner sep=3pt,text=black},
}
\newcommand{\shot}[3][\linewidth]{%
  \begin{tikzpicture}
    \node[anchor=south west,inner sep=0] (img) at (0,0) {\includegraphics[width=#1]{fig/#2}};
    \begin{scope}[x={(img.south east)},y={(img.north west)}]
      #3
    \end{scope}
  \end{tikzpicture}}
% 番号つきの丸（本文用）
\newcommand{\cnum}[1]{\tikz[baseline=(n.base)]\node[num,font=\bfseries\sffamily\scriptsize,minimum size=1.2em](n){#1};}
% 画面の文字（ボタン名など）
\newcommand{\ui}[1]{\textsf{\textbf{「#1」}}}
\newcommand{\key}[1]{\tikz[baseline=(k.base)]\node[draw=gray,rounded corners=2pt,inner sep=1.5pt,font=\sffamily\small,fill=gray!8](k){#1};}

\setlist{itemsep=2pt,topsep=4pt}
\renewcommand{\arraystretch}{1.25}

\title{\includegraphics[width=3.2cm]{../../app/resources/icon/icon_1024.png}\\[1em]
  {\Huge\sffamily\bfseries RawHDR Composer}\\[0.6em]{\LARGE 使用説明書}}
\author{RawHDR Composer 1.1.0}
\date{2026年10月}

\begin{document}
\maketitle

\begin{tcolorbox}[colback=notebg,colframe=notefr,title={\faBookOpen\ この説明書について},fonttitle=\bfseries\sffamily]
RawHDR Composer は、明るさを変えて撮った何枚かの RAW 写真を 1 枚の RAW（DNG）にまとめるアプリです。
この説明書では、初めて使う方でも迷わないよう、画面の写真に\textcolor{accent}{\textbf{赤い矢印・番号・枠}}を描いて、
どこを見て、どこを押せばよいかを順に説明します。

\begin{itemize}
  \item すぐに使いたい方は、\textbf{第\ref{ch:quick}章「基本の使い方」}だけ読めば使えます。
  \item 設定の数値の意味や、変えるときの目安は\textbf{第\ref{ch:settings}章}にまとめています。
  \item うまくいかないときは\textbf{第\ref{ch:faq}章「困ったときは」}を見てください。
\end{itemize}
\end{tcolorbox}

\tableofcontents

% =====================================================================================
\chapter{はじめに}
% =====================================================================================

\section{このアプリでできること}

夕暮れの岩場と月、夜の灯台のように、\textbf{明るい所と暗い所の差がとても大きい場面}は、1 枚の写真では写しきれません。
明るい所に合わせると暗い所が真っ黒につぶれ、暗い所に合わせると明るい所が真っ白に飛んでしまいます。

そこで、カメラの\textbf{AEB（オートブラケット）}で明るさを変えて何枚か撮り、このアプリで 1 枚にまとめます。
まとめた結果は\textbf{RAW のまま}（DNG 形式）なので、Lightroom Classic や Photoshop（Camera Raw）で、
普通の RAW と同じようにホワイトバランス・明るさ・色を自由に調整できます。

\begin{figure}[htbp]
  \centering
  \begin{tikzpicture}
    \node[inner sep=0] (a) {\includegraphics[width=0.48\linewidth]{fig/r_col_single.jpg}};
    \node[inner sep=0,right=0.02\linewidth of a] (b) {\includegraphics[width=0.48\linewidth]{fig/r_col_50_edit.jpg}};
    \node[above=2pt of a,font=\sffamily\small] {1 枚の RAW（真ん中の明るさの 1 枚）};
    \node[above=2pt of b,font=\sffamily\small] {7 枚を合成した DNG を Lightroom で仕上げたもの};
  \end{tikzpicture}
  \caption{左は 1 枚の RAW をそのまま開いたもの。岩は暗く沈み、月は白く飛んでいます。右は 7 枚をこのアプリで合成し、
  Lightroom でハイライト・シャドウなどを少し整えたもの。岩肌と月の模様の両方が見えます。}
\end{figure}

\paragraph{主な特長}
\begin{itemize}
  \item \textbf{RAW のまま合成}: 色補間や色の変換をする前の値で合成するので、後から何でも調整できます。
  \item \textbf{場所ごとに最適な 1 枚を選ぶ}: 各場所で「白飛びしていない中で一番明るい 1 枚」を使います。
        暗い所は明るく撮った 1 枚（ノイズが少ない）、明るい所は暗く撮った 1 枚（白飛びしていない）から取ります。
        切り替わる所は滑らかにつなぎます。
  \item \textbf{Lightroom で仕上げやすく}: 明暗差を縮めておく機能（明暗差の圧縮）があり、Lightroom の普通のスライダーだけで仕上げられます。
  \item \textbf{レンズ補正も使える}: 元の RAW の撮影情報・レンズ情報を引き継ぐので、Lightroom のレンズプロファイル補正がそのまま使えます。
  \item \textbf{どのメーカーのカメラでも}: Canon・Nikon・Sony・富士フイルム・OM SYSTEM・Panasonic・Pentax など、ほとんどのカメラの RAW を読み込めます。
\end{itemize}

\section{合成のしくみ（かんたんな説明）}

\begin{figure}[htbp]
\centering
\begin{tikzpicture}[x=1cm,y=1cm,font=\sffamily\small]
  % 明るさの目盛り
  \shade[left color=black,right color=white] (0,0) rectangle (12,0.5);
  \node[below left] at (0,0) {暗い};
  \node[below right] at (12,0) {明るい};
  \node[below] at (6,-0.05) {場面の明るさ（月は岩より 10 段ほど明るい）};
  % 各フレームが写せる範囲
  \foreach \i/\x/\c/\lab in {1/7.2/blue!70/{暗く撮った 1 枚（1/320 秒）},2/4.6/teal!70/{中くらい（1/10 秒）},3/2.0/magenta!70/{明るく撮った 1 枚（3.2 秒）}}{
    \pgfmathsetmacro{\y}{0.9+0.75*\i}
    \fill[\c,opacity=0.25] (\x,\y) rectangle ++(4.6,0.5);
    \draw[\c,line width=1pt] (\x,\y) rectangle ++(4.6,0.5);
    \node[anchor=east] at (\x-0.1,\y+0.25) {\lab};
    \node[anchor=west,font=\sffamily\scriptsize] at (\x+4.65,\y+0.25) {白飛び};
  }
  \node[anchor=west,text width=11.5cm,align=left] at (0,4.3) {各フレームが\textbf{ちょうどよく写せる範囲}（色の帯）。暗い所は明るく撮った 1 枚、明るい所は暗く撮った 1 枚から取ってつなぎます。};
\end{tikzpicture}
\caption{露出を変えた何枚かで、場面全体の明るさを分担して写す。}
\end{figure}

\section{用意するもの}
\begin{itemize}
  \item \textbf{Mac}: macOS 10.12.6（Sierra）以降。Intel・Apple シリコンのどちらでも動きます。
  \item \textbf{RAW 写真}: 同じ場面を、明るさだけを変えて撮った 2 枚以上（撮り方は第\ref{ch:shoot}章）。
  \item \textbf{現像ソフト}: 書き出した DNG を仕上げるのに、Lightroom Classic・Lightroom・Photoshop（Camera Raw）などを使います。
  \item \textbf{Adobe DNG Converter}（無料・任意）: 無くても動きます。入れておくと、色と明るさが Lightroom で元の RAW を開いたときとそろいます（\ref{sec:dngconv}節）。
\end{itemize}

% =====================================================================================
\chapter{撮影のしかた}\label{ch:shoot}
% =====================================================================================

合成の出来は、撮り方で大きく変わります。次の点を守って撮ってください。

\begin{enumerate}
  \item \textbf{RAW で撮る}（JPEG は使えません）。
  \item \textbf{三脚を使う}。合成する何枚かがずれないようにします。シャッターは 2 秒タイマーかリモコンで切ると、ぶれません。
  \item \textbf{AEB（オートブラケット）を使う}。明るさの間隔は 1〜2 段（±1〜2 EV）、枚数は 3〜7 枚が目安です。
  \item \textbf{絞り・ISO 感度・ピントは固定}。明るさはシャッター速度だけで変えます（絞り優先またはマニュアル。ISO は「オート」にしない）。
  \item \textbf{範囲を確かめる}:
    \begin{itemize}
      \item \textbf{一番暗い 1 枚}で、一番明るい所（月・光源）が\textbf{白飛びしていない}こと。
      \item \textbf{一番明るい 1 枚}で、残したい暗い所（岩・草地）が\textbf{十分に明るく}写っていること。
    \end{itemize}
    カメラのヒストグラムや白飛び警告で確かめ、足りなければ枚数や間隔を増やします。
\end{enumerate}

\begin{tip}[枚数と間隔の目安]
\begin{tabularx}{\linewidth}{lX}
  \toprule
  場面 & 目安 \\ \midrule
  夕景・室内と窓の外 & 2 段おきに 3〜5 枚 \\
  夜景の灯台・街灯 & 2〜2.7 段おきに 5 枚 \\
  月と地上（薄明） & 1.7 段おきに 7 枚（この説明書の作例） \\
  \bottomrule
\end{tabularx}
\end{tip}

\begin{warn}[動くものに注意]
撮っている間に動いたもの（人・車・波・雲・月）は、フレームごとに位置が違うため、つなぎ目で二重に見えたり色がにじんだりすることがあります。
動くものが主役のときは、枚数を減らし、間隔を詰めて手早く撮ってください。
\end{warn}

% =====================================================================================
\chapter{インストールと初めての起動}
% =====================================================================================

\section{インストール}

配布しているディスクイメージ \textsf{RawHDR-Composer-（版）.dmg}（例: \textsf{RawHDR-Composer-1.1.0.dmg}） をダウンロードし、ダブルクリックして開くと、次のウインドウが出ます。

\begin{figure}[htbp]
\centering
\shot[0.62\linewidth]{c_dmg.jpg}{
  \node[num] at (0.15,{1-0.32}) {1};
  \node[num] at (0.67,{1-0.32}) {2};
  \node[num] at (0.15,{1-0.71}) {3};
  \node[num] at (0.67,{1-0.71}) {4};
  \draw[arr,dashed] (0.31,{1-0.32}) -- (0.62,{1-0.32});
}
\caption{ディスクイメージの中身}\label{fig:dmg}
\end{figure}

\begin{enumerate}[label=\protect\cnum{\arabic*}]
  \item \textsf{RawHDR Composer.app}: アプリ本体。
  \item \ui{アプリケーション}: 「アプリケーション」フォルダへの入り口（手で入れるとき、ここへ\cnum{1}をドラッグします）。
  \item \textsf{かんたんインストーラ.scpt}: アプリを入れる作業を自動で行います（おすすめ）。
  \item \textsf{使用説明書.pdf}: この説明書です。
\end{enumerate}

\subsection{かんたんインストーラで入れる（おすすめ）}

かんたんインストーラは、次の 3 つを自動で行います。
\begin{itemize}
  \item アプリを\ui{アプリケーション}フォルダへコピーする
  \item 「開発元を確認できない」で開けない状態を、\textbf{このアプリについてだけ}解除する
  \item アプリを初めて起動する
\end{itemize}

\begin{enumerate}
  \item ウインドウの\textsf{かんたんインストーラ.scpt}（図\ref{fig:dmg}の\cnum{3}）をダブルクリックします。\ui{スクリプトエディタ}が開きます。
  \item スクリプトエディタの右上にある\textbf{▶（実行）ボタン}を押します。
  \item 確認の画面で\ui{インストール}を押します。Mac のログインパスワードを求められたら入力します。
        すでにインストールしてあるときは、置き換えてよいかを聞かれます。
  \item アプリが起動し、「インストールが完了し…」と出たら\ui{OK}を押して終わりです。スクリプトエディタは閉じて構いません（保存はしません）。
  \item Finder のサイドバーの\ui{RawHDR Composer（版）}の横にある取り出しボタン（\faEject）で、ディスクイメージを取り出します。
\end{enumerate}

\begin{figure}[htbp]
\centering
\shot[0.8\linewidth]{c_script_editor.jpg}{
  \draw[hl] (0.785,{1-0.02}) rectangle (0.835,{1-0.08});
  \node[lbl,anchor=north east] (r) at (0.78,{1-0.16}) {\textbf{▶（実行）を押す}};
  \draw[arr] (r.north east) -- (0.80,{1-0.085});
  \node[lbl,anchor=north west,font=\sffamily\scriptsize] at (0.02,{1-0.80}) {中身（AppleScript）は書き換えなくて構いません};
}
\caption{かんたんインストーラをスクリプトエディタで開いたところ}
\end{figure}

\begin{memo}[かんたんインストーラが行う「解除」について]
インターネットから入手したファイルには、macOS が「検疫情報」という印を付け、初めて開くときに確かめます。
このアプリは Apple の公証を受けていないため、その確かめで止められることがあります。
かんたんインストーラは、\ui{アプリケーション}フォルダに入れた RawHDR Composer \textbf{だけ}からこの印を外します。
macOS の Gatekeeper そのものを止めたり、ほかのアプリに何かをしたりはしません。
途中で「スクリプトエディタが〜を変更しようとしています」などの確認が出たときは、許可してください。
\end{memo}

\subsection{手で入れる}
\begin{enumerate}
  \item ウインドウの\textsf{RawHDR Composer.app}（図\ref{fig:dmg}の\cnum{1}）を、\ui{アプリケーション}（\cnum{2}）へドラッグします。
  \item \ui{アプリケーション}フォルダの \textsf{RawHDR Composer} をダブルクリックして起動します。
\end{enumerate}

\begin{warn}[「開発元を確認できない」と出て開けないとき]
このアプリは Apple の公証を受けていないため、初めて開くときに macOS が止めることがあります。かんたんインストーラを使うか、次のようにして開いてください（2 回目からは普通に開けます）。
\begin{itemize}
  \item \textbf{macOS 14 以前}: Finder でアプリを \key{control} を押しながらクリック（または右クリック）し、\ui{開く}を選ぶ。出てきた確認で\ui{開く}を押す。
  \item \textbf{macOS 15 以降}: 一度ダブルクリックして警告を閉じたあと、\ui{システム設定}→\ui{プライバシーとセキュリティ}を開き、下の方にある「RawHDR Composer」の\ui{このまま開く}を押す。
\end{itemize}
\end{warn}

\section{Adobe DNG Converter の案内}\label{sec:dngconv}

Adobe DNG Converter（Adobe の無料アプリ）が入っていないと、起動したときに次の案内が出ます。

\begin{figure}[htbp]
\centering
\shot[0.92\linewidth]{c_dng_prompt.jpg}{
  \draw[hl] (0.17,{1-0.165}) rectangle (0.965,{1-0.265});
  \node[lbl,anchor=west] (t1) at (1.03,{1-0.215}) {\textbf{無くても}\\\textbf{すべて動きます}};
  \draw[arr] (t1.west) -- (0.965,{1-0.215});
  \node[num] at (0.32,{1-1.0}) {1};
  \node[num] at (0.50,{1-1.0}) {2};
  \node[num] at (0.79,{1-1.0}) {3};
  \node[num] at (0.155,{1-0.80}) {4};
}
\caption{Adobe DNG Converter の入手の案内}
\end{figure}

\begin{enumerate}[label=\protect\cnum{\arabic*}]
  \item \ui{あとで}: 今は入れずに使い始めます。アプリはこのまま\textbf{すべての機能が使えます}。
  \item \ui{Adobe のページを開く}: Adobe の DNG Converter の説明ページをブラウザで開きます。
  \item \ui{ダウンロードしてインストール…}: Adobe の公式サイトから最新版（約 1.8 GB）をダウンロードし、インストーラーを開きます。
        インストーラーの指示に従って、使用許諾を確認しながらインストールしてください。インストールが済むと、このアプリが自動で見つけて使います（再起動は不要）。
  \item \ui{起動時にこの案内を表示しない}: チェックすると、次からこの案内は出ません。
\end{enumerate}

\begin{memo}[入れると何が変わる？]
入れておくと、書き出す DNG の\textbf{色・明るさ・カメラプロファイル・ノイズの情報}が、Lightroom で元の RAW を開いたときとそろいます。
入れていないときは、色と明るさが元の RAW とわずかに違います（仕上げには困らない程度です）。
後から入れたくなったら、メニュー\ui{RawHDR Composer}→\ui{Adobe DNG Converter…}でこの案内を開けます。
最新版の DNG Converter は macOS 13 以降で動きます。それより古い macOS では、案内の\ui{Adobe のページを開く}から対応する版を確かめてください。
\end{memo}

% =====================================================================================
\chapter{画面の見方}
% =====================================================================================

画面は、\textbf{左にフレームの一覧}、\textbf{中央にプレビュー}、\textbf{右に設定と書き出し}が並んでいます。

\begin{figure}[htbp]
\centering
\shot{04_merged.jpg}{
  \draw[hl] (0.003,{1-0.012}) rectangle (0.232,{1-0.89});
  \node[num] at (0.20,{1-0.30}) {1};
  \draw[hl] (0.003,{1-0.893}) rectangle (0.155,{1-0.925});
  \node[num] at (0.172,{1-0.909}) {2};
  \draw[hl] (0.235,{1-0.008}) rectangle (0.552,{1-0.038});
  \node[num] at (0.57,{1-0.023}) {3};
  \draw[hl] (0.235,{1-0.040}) rectangle (0.39,{1-0.066});
  \node[num] at (0.39,{1-0.09}) {4};
  \draw[hl] (0.392,{1-0.040}) rectangle (0.48,{1-0.066});
  \node[num] at (0.47,{1-0.09}) {5};
  \draw[hl] (0.483,{1-0.040}) rectangle (0.612,{1-0.066});
  \node[num] at (0.63,{1-0.053}) {6};
  \node[num] at (0.53,{1-0.50}) {7};
  \draw[hl] (0.82,{1-0.010}) rectangle (0.996,{1-0.115});
  \node[num] at (0.80,{1-0.07}) {8};
  \draw[hl] (0.82,{1-0.150}) rectangle (0.988,{1-0.627});
  \node[num] at (0.80,{1-0.39}) {9};
  \draw[hl] (0.82,{1-0.630}) rectangle (0.988,{1-0.658});
  \node[num] at (0.80,{1-0.644}) {10};
  \draw[hl] (0.82,{1-0.668}) rectangle (0.988,{1-0.812});
  \node[num] at (0.80,{1-0.74}) {11};
  \draw[hl] (0.82,{1-0.825}) rectangle (0.988,{1-0.972});
  \node[num] at (0.80,{1-0.90}) {12};
  \draw[hl] (0.003,{1-0.965}) rectangle (0.40,{1-0.995});
  \node[num] at (0.42,{1-0.98}) {13};
}
\caption{合成した後の画面}\label{fig:overview}
\end{figure}

\begin{tabularx}{\linewidth}{>{\centering\arraybackslash}p{2.2em}l X}
\toprule
番号 & 名前 & 役割 \\ \midrule
\cnum{1} & フレームの一覧 & 読み込んだ RAW の一覧（暗い順）。選ぶとプレビューにそのフレームが出る。番号の色は「由来マップ」の色 \\
\cnum{2} & 一覧の操作 & \ui{追加…}（RAW を選んで加える）・\ui{取り除く}（選んだフレームを外す）・\ui{すべて取り除く} \\
\cnum{3} & 表示の切り替え & プレビューに何を出すか（\ref{sec:views}節） \\
\cnum{4} & 表示の明るさ & プレビューだけを明るく・暗く（$-6$〜$+6$ 段）。書き出す DNG には影響しない \\
\cnum{5} & ハイライトを見やすく & プレビューで、明るすぎる所の模様も見えるようにする（表示だけ） \\
\cnum{6} & 拡大縮小 & \ui{−}\ui{＋}\ui{全体}\ui{100\%}と、今の倍率 \\
\cnum{7} & プレビュー & ドラッグで動かし、ホイール・ピンチで拡大縮小 \\
\cnum{8} & 基準フレーム & DNG の明るさ・撮影情報の基準にするフレーム \\
\cnum{9} & 合成の設定 & つなぎ方・ノイズを減らす平均・ゴースト対策・明暗差の圧縮など（第\ref{ch:settings}章） \\
\cnum{10} & HDR 結合のボタン & 最初の合成を始める（\ui{HDR結合開始}）。その後は設定を変えると自動で合成し直す（\ui{再結合中…}→\ui{結合済み}） \\
\cnum{11} & 書き出し & DNG の形式と、\ui{DNG を書き出す…}のボタン \\
\cnum{12} & 解析の結果 & 露出の測り方・合成の内訳など（第\ref{ch:info}章）。右の欄は、ウインドウが低いときは縦にスクロールできる \\
\cnum{13} & 状態の表示 & いま何をしているか、次に何をすればよいか \\
\bottomrule
\end{tabularx}

% =====================================================================================
\chapter{基本の使い方}\label{ch:quick}
% =====================================================================================

RAW を読み込み、合成し、DNG を書き出して、Lightroom で仕上げるまでの流れです。

% ---------------------------------------------------------------------------------
\begin{step}{1}{RAW を読み込む}
同じ場面を明るさを変えて撮った RAW を、\textbf{まとめてウインドウへドラッグ＆ドロップ}します（フォルダごとでも構いません）。
左下の\ui{追加…}を押して選ぶこともできます。

\centering
\shot[0.95\linewidth]{01_empty.jpg}{
  \node[lbl,anchor=north] (d) at (0.53,{1-0.62}) {\textbf{ここへ RAW をドラッグ＆ドロップ}\\（何枚でも、フォルダごとでも可）};
  \node[inner sep=0] (f) at (0.53,{1-0.85}) {\Huge\textcolor{accent}{\faImages}};
  \draw[arr] (f.north) -- ++(0,0.08);
  \draw[hl] (0.003,{1-0.893}) rectangle (0.045,{1-0.925});
  \node[lbl,anchor=west] (a) at (0.06,{1-0.80}) {または\ui{追加…}で選ぶ};
  \draw[arr] (a.south west) -- (0.03,{1-0.89});
}
\end{step}

\begin{figure}[htbp]
\centering
\shot[0.75\linewidth]{c_open.jpg}{
  \draw[hl] (0.245,{1-0.215}) rectangle (0.99,{1-0.42});
  \node[lbl,anchor=west] (s) at (1.02,{1-0.30}) {\key{shift}を押しながら\\クリックして\\まとめて選ぶ};
  \draw[arr] (s.west) -- (0.99,{1-0.30});
  \node[num] at (0.885,{1-0.91}) {\faCheck};
  \node[lbl,anchor=west] (o) at (1.02,{1-0.91}) {\ui{開く}};
  \draw[arr] (o.west) -- (0.91,{1-0.91});
}
\caption{\ui{追加…}を押したときのファイルの選択（左の欄はぼかしています）}
\end{figure}

読み込むと、一覧に\textbf{暗い順}に並び、プレビューに 1 枚目のフレームが表示されます。
このときはまだ合成していません。

% ---------------------------------------------------------------------------------
\begin{step}{2}{フレームを確かめる}
一覧のフレームを\textbf{クリックすると、そのフレームがプレビューに出ます}。
一番暗い 1 枚で月や光源が白飛びしていないか、一番明るい 1 枚で暗い所が十分に写っているかを確かめましょう。

\centering
\shot[0.95\linewidth]{02_loaded.jpg}{
  \draw[hl] (0.005,{1-0.205}) rectangle (0.232,{1-0.236});
  \node[lbl,anchor=north west] (c) at (0.02,{1-0.30}) {\textbf{クリック}すると\\そのフレームを表示};
  \draw[arr] (c.north) -- (0.10,{1-0.24});
  \draw[arr] (0.233,{1-0.22}) to[bend left=20] (0.40,{1-0.30});
  \draw[hl] (0.360,{1-0.008}) rectangle (0.418,{1-0.038});
  \node[lbl,anchor=north] (m) at (0.40,{1-0.10}) {表示が自動で\ui{選んだフレーム}になる};
  \draw[arr] (m.north) -- (0.389,{1-0.04});
  \node[lbl,anchor=north] (r) at (0.12,{1-0.45}) {一覧の\textbf{相対EV}: 一番暗い 1 枚\\より何段明るいか（実測値）};
}
\end{step}

\begin{memo}[一覧の列の意味]
\textbf{\#}: 番号（色は由来マップの色）／\textbf{ファイル}: ファイル名／\textbf{シャッター}: シャッター速度／
\textbf{相対EV}: 一番暗い 1 枚より何段明るいか（写真の中身から測った値）／\textbf{ずれ}: 位置合わせで動かした量（横,縦 の画素。回転は解析の結果に出る）。
\textbf{太字}の行が基準フレームです。行にマウスを重ねると、F 値・ISO・レンズなどが出ます。
\end{memo}

% ---------------------------------------------------------------------------------
\begin{step}{3}{「HDR結合開始」を押す}
右の欄の\ui{HDR結合開始}を押すと合成が始まります（\key{⌘}\key{R} でも可）。
数十秒ほどで終わり、プレビューが合成結果に切り替わります。

\centering
\shot[0.95\linewidth]{03_before_merge.jpg}{
  \draw[hl] (0.82,{1-0.630}) rectangle (0.988,{1-0.658});
  \node[lbl,anchor=east] (b) at (0.79,{1-0.56}) {\textbf{押す}と合成が始まる\\（青いボタンは Return キーでも押せる）};
  \draw[arr] (b.east) -- (0.85,{1-0.635});
  \draw[hl] (0.82,{1-0.771}) rectangle (0.988,{1-0.798});
  \node[lbl,anchor=east] (e) at (0.79,{1-0.84}) {合成するまでは\\書き出せない（灰色）};
  \draw[arr] (e.east) -- (0.85,{1-0.79});
}
\end{step}

% ---------------------------------------------------------------------------------
\begin{step}{4}{合成結果を確かめる}
合成が終わると、プレビューに合成結果が出ます。拡大して、つなぎ目や月・光源の周りを確かめましょう。
どのフレームがどこに使われたかは\ui{由来マップ}で見られます（\ref{sec:views}節）。

\centering
\shot[0.95\linewidth]{04_merged.jpg}{
  \draw[hl] (0.235,{1-0.008}) rectangle (0.272,{1-0.038});
  \node[lbl,anchor=north west] (v) at (0.25,{1-0.10}) {\ui{合成結果}の表示};
  \draw[arr] (v.north west) -- (0.255,{1-0.04});
  \draw[hl] (0.483,{1-0.040}) rectangle (0.612,{1-0.066});
  \node[lbl,anchor=north] (z) at (0.60,{1-0.13}) {拡大して細部を確かめる\\（ドラッグで移動）};
  \draw[arr] (z.north) -- (0.55,{1-0.068});
  \draw[hl] (0.82,{1-0.630}) rectangle (0.988,{1-0.658});
  \node[lbl,anchor=east] (k) at (0.79,{1-0.58}) {\ui{結合済み}になる};
  \draw[arr] (k.east) -- (0.85,{1-0.638});
  \draw[hl] (0.003,{1-0.965}) rectangle (0.30,{1-0.995});
}
\end{step}

% ---------------------------------------------------------------------------------
\begin{step}{5}{DNG を書き出す}
右の欄の\ui{DNG を書き出す…}を押し（\key{⌘}\key{E} でも可）、保存する場所と名前を決めて\ui{保存}を押します。
名前と場所は、はじめは基準フレームと同じフォルダの「（基準フレームの名前）\_HDR.dng」になっています。

\centering
\shot[0.44\linewidth]{c_right_merged.jpg}{
  \draw[hl] (0.03,{1-0.952}) rectangle (0.97,{1-0.984});
  \node[lbl,anchor=east] (x) at (-0.05,{1-0.968}) {\textbf{押して保存}};
  \draw[arr] (x.east) -- (0.03,{1-0.968});
  \draw[hl] (0.03,{1-0.845}) rectangle (0.97,{1-0.871});
  \node[lbl,anchor=east] (f) at (-0.05,{1-0.85}) {形式は\\\ui{自動（おすすめ）}\\のままでよい};
  \draw[arr] (f.east) -- (0.03,{1-0.858});
}
\end{step}

\begin{figure}[htbp]
\centering
\shot[0.6\linewidth]{c_save.jpg}{
  \draw[hl] (0.22,{1-0.07}) rectangle (0.56,{1-0.22});
  \node[lbl,anchor=east] (n) at (-0.03,{1-0.15}) {名前};
  \draw[arr] (n.east) -- (0.22,{1-0.15});
  \draw[hl] (0.22,{1-0.54}) rectangle (0.62,{1-0.68});
  \node[lbl,anchor=east] (p) at (-0.03,{1-0.61}) {保存する場所};
  \draw[arr] (p.east) -- (0.22,{1-0.61});
  \node[lbl,anchor=west] (s) at (1.03,{1-0.84}) {\ui{保存}};
  \draw[arr] (s.west) -- (0.92,{1-0.84});
}
\caption{保存の画面}
\end{figure}

% ---------------------------------------------------------------------------------
\begin{step}{6}{Lightroom で仕上げる}
書き出した DNG を Lightroom Classic に読み込みます（または Photoshop で開くと Camera Raw が開きます）。
普通の RAW と同じように仕上げられます。おすすめの手順は第\ref{ch:lr}章で説明します。
\end{step}

\begin{tip}[設定を変えると、自動で合成し直す]
一度合成した後に右の欄の設定を変えると、少し待ってから\textbf{自動で合成し直します}（スライダーを動かしている間は待ちます）。
合成し直している間はボタンが\ui{再結合中…}になり、書き出しのボタンは押せません（古い結果をうっかり書き出さないため）。
\ui{結合済み}に戻ったら書き出せます。最初の 1 回だけは、設定を確かめてから\ui{HDR結合開始}を押します。

\centering
\shot[0.42\linewidth]{c_right_changed.jpg}{
  \draw[hl] (0.03,{1-0.773}) rectangle (0.97,{1-0.800});
  \node[lbl,anchor=east] (r) at (-0.05,{1-0.786}) {\ui{再結合中…}};
  \draw[arr] (r.east) -- (0.03,{1-0.786});
  \draw[hl] (0.03,{1-0.535}) rectangle (0.97,{1-0.582});
  \node[lbl,anchor=east] (c) at (-0.05,{1-0.558}) {設定を変えた};
  \draw[arr] (c.east) -- (0.03,{1-0.558});
}
\end{tip}

% =====================================================================================
\chapter{プレビューの使い方}
% =====================================================================================

\section{拡大・縮小と移動}

\begin{figure}[htbp]
\centering
\shot{c_toolbar.jpg}{
  \foreach \x/\w/\n in {0.683/0.020/1,0.737/0.020/2,0.804/0.035/3,0.888/0.044/4,0.955/0.020/5}{
    \draw[hl] (\x-\w,{1-0.55}) rectangle (\x+\w,{1-0.93});
    \node[num] at (\x,{1-1.15}) {\n};}
  \draw[hl] (0.008,{1-0.55}) rectangle (0.405,{1-0.93});
  \node[num] at (0.21,{1-1.15}) {6};
  \draw[hl] (0.415,{1-0.55}) rectangle (0.64,{1-0.93});
  \node[num] at (0.52,{1-1.15}) {7};
}
\caption{プレビューの上の操作}
\end{figure}

\begin{enumerate}[label=\protect\cnum{\arabic*}]
  \item \ui{−}: 縮小（\key{⌘}\key{−}）
  \item \ui{＋}: 拡大（\key{⌘}\key{＋}）
  \item \ui{全体}: 画像全体が収まる大きさにする（\key{⌘}\key{0}）
  \item \ui{100\%}: プレビューの 1 画素を画面の 1 画素にする（\key{⌘}\key{1}）
  \item 今の倍率
  \item 表示の明るさ: プレビューだけを明るく・暗くする（$-6$〜$+6$ 段、0.5 段刻み）。暗部を確かめるときに上げます。\textbf{書き出す DNG には影響しません。}
  \item \ui{ハイライトを見やすく}: 明るすぎる所（月・光源）の模様も見えるように、プレビューの明るさの分布を縮めます（表示だけ）。
        明暗差の圧縮（\ref{sec:compress}節）を使っているときは、Lightroom で開いたときと同じ見え方にするため働きません。
\end{enumerate}

\begin{tip}[プレビューの動かし方]
\begin{itemize}
  \item \textbf{移動}: 画像の上で\textbf{ドラッグ}します（カーソルが手のひら \faHandPaper[regular] になります）。
  \item \textbf{拡大縮小}: マウスの\textbf{ホイール}、トラックパッドの\textbf{ピンチ}、または上のボタン。ホイールではカーソルの位置を中心に拡大縮小します。
\end{itemize}
\end{tip}

\section{表示の切り替え}\label{sec:views}

プレビューの上の 7 つのボタンで、表示する内容を切り替えます。合成の様子を確かめるのに使います。

\subsection{合成結果}
合成した結果です。書き出す DNG を Lightroom・Camera Raw で\textbf{何も調整せずに開いたときに近い見え方}です（図\ref{fig:overview}）。
書き出す DNG と同じ画素（色補間して、明暗差の圧縮の倍率を画素ごとに配り直したもの）を、DNG を開いたときと同じ明るさで、
Camera Raw の既定の見え方に近いトーンカーブで描いています。ただし、Lightroom で選んだプロファイル（カメラのポートレートなど）による色の違いは再現しません。

\subsection{由来マップ・重ねて表示}
\textbf{由来マップ}は、それぞれの場所で\textbf{どのフレームを使ったか}を色で示します。色は一覧の番号の色と同じです。
\textbf{重ねて表示}は、合成結果に由来マップを薄く重ねたものです。

\begin{figure}[htbp]
\centering
\shot[0.9\linewidth]{11_lh_sourcemap.jpg}{
  \draw[hl] (0.003,{1-0.06}) rectangle (0.05,{1-0.145});
  \node[lbl,anchor=north west] (l) at (0.02,{1-0.25}) {番号の色\\＝由来マップの色};
  \draw[arr] (l.north) -- (0.03,{1-0.15});
  \node[lbl,anchor=west] (o) at (0.30,{1-0.30}) {橙＝5 番（一番明るく撮った 1 枚）\\：空・草地};
  \node[lbl,anchor=west] (g) at (0.60,{1-0.62}) {緑・黄・青＝暗く撮った\\フレーム：灯台・灯室};
  \draw[arr] (g.west) -- (0.565,{1-0.66});
  \draw[hl] (0.316,{1-0.008}) rectangle (0.359,{1-0.038});
}
\caption{由来マップ（灯台の作例）。暗い所は明るく撮った 1 枚、灯台や灯室は暗く撮ったフレームから取っています。}
\end{figure}

\begin{figure}[htbp]
\centering
\shot[0.9\linewidth]{12_lh_overlay.jpg}{
  \draw[hl] (0.272,{1-0.008}) rectangle (0.315,{1-0.038});
  \node[lbl,anchor=north] (o) at (0.30,{1-0.12}) {\ui{重ねて表示}};
  \draw[arr] (o.north) -- (0.294,{1-0.04});
}
\caption{重ねて表示。合成結果に由来マップの色を薄く重ねます。}
\end{figure}

\subsection{選んだフレーム}
一覧で選んだフレームを 1 枚だけ表示します。一覧をクリックすると自動でこの表示になります。

\subsection{位置の確認}
一覧で選んだフレームと基準フレームを、明るさをそろえて比べます。
\textbf{位置が合っていない所が赤く}、\textbf{比べられない所（どちらかが白飛び・暗すぎる）が青く}、それ以外は灰色で表示されます。
位置合わせ（\ref{sec:align}節）を確かめるときに使います。

\begin{figure}[htbp]
\centering
\shot[0.8\linewidth]{09_align.jpg}{
  \draw[hl] (0.418,{1-0.008}) rectangle (0.463,{1-0.038});
  \node[lbl,anchor=north west,font=\sffamily\scriptsize] (a) at (0.35,{1-0.12}) {\ui{位置の確認}};
  \draw[arr] (a.north) -- (0.44,{1-0.04});
  \node[lbl,anchor=west,font=\sffamily\scriptsize] (r) at (0.50,{1-0.55}) {赤い所＝ずれている、\\または明るさが違う};
  \draw[arr] (r.west) -- (0.45,{1-0.68});
  \node[lbl,anchor=west,font=\sffamily\scriptsize] (b) at (0.26,{1-0.30}) {青い所＝比べられない\\（白飛び・暗すぎ）};
}
\caption{位置の確認（灯台の作例）。三脚で撮ったのでずれはなく、赤いのは光の当たり方の違う灯台だけです。}
\end{figure}

\subsection{前回との違い}
設定を変えて合成し直したとき、\textbf{前の設定の結果から変わった所を赤く}表示します（1\% の違いで薄く、10\% 以上で濃く）。
設定がどこに効いたのかを確かめるのに使います。下の状態の表示に、変わった面積の割合も出ます。

\begin{figure}[htbp]
\centering
\shot[0.8\linewidth]{08_diff.jpg}{
  \draw[hl] (0.463,{1-0.008}) rectangle (0.513,{1-0.038});
  \node[lbl,anchor=north west,font=\sffamily\scriptsize] (a) at (0.52,{1-0.10}) {\ui{前回との違い}};
  \draw[arr] (a.north west) -- (0.50,{1-0.04});
  \draw[hl] (0.003,{1-0.965}) rectangle (0.80,{1-0.995});
  \node[lbl,anchor=south,font=\sffamily\scriptsize] (s) at (0.40,{1-0.86}) {変わった面積の割合};
  \draw[arr] (s.south) -- (0.40,{1-0.962});
}
\caption{明暗差の圧縮を 50\% から 80\% に変えて合成し直したときの「前回との違い」。空と海が大きく変わりました。}
\end{figure}

\subsection{動いた所}
撮影の間に動いた物（人・車・枝・波・雲・光の筋など）を見つけて、\textbf{ゴースト対策（\ref{sec:deghost}節）で重みを変えた所を赤く}表示します。
そこでは、基準フレーム（暗い所ではそれより明るいフレーム）と食い違うフレームを使わないので、動いた物が二重や半透明に写りません。
解析の結果に、見つかった所の数と面積も出ます。

\begin{figure}[htbp]
\centering
\shot[0.8\linewidth]{18_lh_ghost.jpg}{
  \draw[hl] (0.513,{1-0.008}) rectangle (0.551,{1-0.038});
  \node[lbl,anchor=north west,font=\sffamily\scriptsize] (a) at (0.57,{1-0.10}) {\ui{動いた所}};
  \draw[arr] (a.north west) -- (0.545,{1-0.04});
  \node[lbl,anchor=west,font=\sffamily\scriptsize] (b) at (0.66,{1-0.30}) {回転する灯台の光の筋};
  \draw[arr] (b.west) -- (0.585,{1-0.30});
  \node[lbl,anchor=north,font=\sffamily\scriptsize] (g) at (0.30,{1-0.82}) {風で揺れる草・木};
  \draw[arr] (g.north) -- (0.33,{1-0.72});
}
\caption{動いた所（灯台の作例）。灯台の光の筋、風で揺れる草や木、光の当たり方が変わった塔が見つかりました。}
\end{figure}

% =====================================================================================
\chapter{設定項目の説明}\label{ch:settings}
% =====================================================================================

右の欄の設定を、上から順に説明します。\textbf{はじめはすべて既定値のままで構いません。}
仕上がりに気になる所があったときに、この章の「変えるときの目安」を参考に調整してください。

\begin{figure}[htbp]
\centering
\shot[0.55\linewidth]{c_right_merged.jpg}{
  \node[num] at (-0.05,{1-0.045}) {A};
  \node[num] at (-0.05,{1-0.115}) {B};
  \node[num] at (-0.05,{1-0.265}) {C};
  \node[num] at (-0.05,{1-0.324}) {D};
  \node[num] at (-0.05,{1-0.381}) {E};
  \node[num] at (-0.05,{1-0.427}) {F};
  \node[num] at (-0.05,{1-0.484}) {G};
  \draw[thick,accent] (-0.015,{1-0.445}) -- (-0.015,{1-0.525});
  \node[num] at (-0.05,{1-0.634}) {H};
  \draw[thick,accent] (-0.015,{1-0.535}) -- (-0.015,{1-0.730});
  \node[num] at (-0.05,{1-0.751}) {I};
  \node[num] at (-0.05,{1-0.786}) {J};
  \node[num] at (-0.05,{1-0.858}) {K};
  \node[num] at (-0.05,{1-0.935}) {L};
}
\caption{右の欄の設定}\label{fig:right}
\end{figure}

\begin{center}
\small
\begin{tabularx}{\linewidth}{c >{\raggedright\arraybackslash}p{11\zw} >{\raggedright\arraybackslash}p{6.5\zw} >{\raggedright\arraybackslash}p{8.5\zw} X}
\toprule
 & 設定 & 既定値 & 範囲 & ひとことで \\ \midrule
A & 基準フレーム & 自動 & 一覧のフレーム & DNG の明るさ・撮影情報の基準 \\
B & 位置合わせ & 自動 & しない／自動／手動 & フレームどうしのずれ・回転を直す \\
C & 切り替えを始める明るさ & 0.55 & 0.30〜0.90 & 暗いフレームへ切り替え始める明るさ \\
D & 飽和とみなす閾値 & 0.92 & 0.80〜0.99 & どこから「白飛び」とみなすか \\
E & 重みのぼかし & 16 px & 4〜256 px & 切り替わりのなめらかさ \\
F & 平均してノイズを減らす & オン & オン／オフ & 複数のフレームを平均して暗部のノイズを減らす \\
G & ゴーストを取り除く・感度 & オン・50\% & オン／オフ・0〜100\% & 動いた物の二重写りを防ぐ \\
H & 明暗差の圧縮 & 50\%（自動はオフ）、明るい所・暗い所 100\% & 0〜100\%、各 0〜200\% & 明暗差を縮めて仕上げやすく \\
I & 色の縁取りを補正 & オン & オン／オフ & レンズの色ずれを直す \\
J & HDR 結合のボタン & --- & --- & 合成を始める \\
K & 書き出しの形式 & 自動 & 自動／CFA／LinearRaw & DNG の中身の形 \\
L & レンズ補正を有効にして開く & オン & オン／オフ & DNG を開いたときにレンズ補正を入れておく \\
\bottomrule
\end{tabularx}
\end{center}

% ---------------------------------------------------------------------------------
\section{A. 基準フレーム}
書き出す DNG を、\textbf{どのフレームと同じ明るさ・撮影情報（日時・シャッター速度・レンズなど）で開くか}を決めます。
\begin{itemize}
  \item \textbf{既定}: \ui{自動（露光量が中央のもの）}。真ん中の明るさのフレームが選ばれます。
  \item 合成の中身（どこにどのフレームを使うか）は、基準フレームを変えてもほとんど変わりません。
  \item \textbf{変えるとき}: DNG を開いたときの明るさを変えたい、または撮影情報を特定の 1 枚にそろえたいとき。
        明るさは Lightroom の露光量でも変えられるので、普通は自動のままで構いません。
\end{itemize}

% ---------------------------------------------------------------------------------
\section{B. 位置合わせ}\label{sec:align}
フレームどうしが少しずれているとき（三脚が動いた、手持ちで撮った）に、ずれ（上下左右の移動・回転・1 画素より細かいずれ）を直します。

\noindent\begin{tabularx}{\linewidth}{l X}
\toprule
選択肢 & 内容 \\ \midrule
しない（三脚） & ずれを直さない \\
自動（既定） & 写真の中身から各フレームのずれ・回転を測って直す。三脚で撮ってずれが無ければ、何も動かさない。一覧の\textbf{ずれ}の列に動かした量（横,縦 の画素）が、解析の結果に回転の角度が出る \\
手動（自動の結果を微調整） & 自動の結果から始めて、矢印のボタンで 1 段（2 画素）ずつ動かす（自動で求めた回転はそのまま） \\
\bottomrule
\end{tabularx}

\begin{figure}[htbp]
\centering
\shot[0.5\linewidth]{c_right_align.jpg}{
  \draw[hl] (0.02,{1-0.306}) rectangle (0.98,{1-0.389});
  \node[lbl,anchor=east,font=\sffamily\scriptsize] (m) at (-0.04,{1-0.347}) {\ui{手動}を選ぶ};
  \draw[arr] (m.east) -- (0.02,{1-0.347});
  \draw[hl] (0.02,{1-0.405}) rectangle (0.98,{1-0.494});
  \node[lbl,anchor=east,font=\sffamily\scriptsize] (n) at (-0.04,{1-0.56}) {一覧で選んだフレームを\\矢印で動かす。\ui{0}で元に戻す};
  \draw[arr] (n.east) -- (0.02,{1-0.45});
}
\caption{手動の位置合わせ}
\end{figure}

\paragraph{手動で合わせる手順}
\begin{enumerate}
  \item \ui{位置合わせ}で\ui{手動（自動の結果を微調整）}を選ぶ（まず自動で合わせます）。
  \item 一覧で、動かしたいフレームを選ぶ（基準フレーム以外）。
  \item プレビューが\ui{位置の確認}になるので、\textbf{赤い輪郭が消える}方向へ矢印のボタンで動かす。
  \item 合わせ終わると自動で合成し直します（最初の合成の前なら\ui{HDR結合開始}を押す）。
\end{enumerate}

\begin{memo}[位置合わせのしくみ]
\begin{itemize}
  \item まず 2 画素単位の上下左右の移動でずれを探し、続けて回転・わずかな拡大縮小・1 画素より細かいずれを詰めます。
  \item 2 画素単位の移動で足りるフレーム（三脚で少し動いた程度）は、画素をそのまま動かします（画質は変わりません）。
  \item 足りないフレーム（手持ちで回ったなど）は、いったん色補間してから回転・移動し、元の RAW の色の並びに戻します。そのフレームから取った所は、わずかに柔らかくなります。
  \item 基準フレームは動かしません。動かした分の端は、DNG の切り抜きで隠れます。
\end{itemize}
\end{memo}

\begin{warn}
三脚で撮っても、手ぶれ補正やピントのわずかな動きで 1 画素ほどずれることがあり、\ui{自動}はそれも直します。
ずれが 1 画素に満たないのに色補間して動かしたフレームは、そこから取った所がわずかに柔らかくなります。気になるときは\ui{しない（三脚）}にしてください。
手持ちで撮った写真での確かめは、まだ十分ではありません。うまく合わないときは\ui{位置の確認}で確かめ、\ui{手動}で直してください。
\end{warn}

% ---------------------------------------------------------------------------------
\section{C. 切り替えを始める明るさ}\label{sec:ramp}
明るく撮ったフレームの値が「白飛びの手前」のどのくらいまで来たら、暗く撮ったフレームへ\textbf{切り替え始めるか}を決めます。
値は「白飛びとみなす明るさ（D）」に対する割合です。

\begin{center}
\begin{tikzpicture}[x=11cm,y=1cm,font=\sffamily\small]
  \draw[line width=1pt] (0,0) -- (1,0);
  \foreach \x/\l in {0.3/0.30,0.55/0.55,0.9/0.90} {\draw (\x,0.1) -- (\x,-0.1) node[below] {\l};}
  \fill[accent] (0.55,0) circle (0.08cm);
  \node[above,accent] at (0.55,0.15) {既定};
  \node[anchor=north west,text width=4.6cm,align=left] at (0,-0.6) {\textbf{小さくする}と\\・切り替わりがゆるやかで継ぎ目が目立ちにくい\\・暗い（ノイズの多い）フレームを使う範囲が広がる};
  \node[anchor=north east,text width=4.6cm,align=left] at (1,-0.6) {\textbf{大きくする}と\\・明るい（ノイズの少ない）フレームを白飛び近くまで使う\\・切り替わりが急になり、段差が出やすい};
\end{tikzpicture}
\end{center}

\begin{tip}[変えるときの目安]
\begin{itemize}
  \item \textbf{ふだん}: 0.55（既定）。
  \item \textbf{空のグラデーションなどに段差・継ぎ目が見える}: 0.40〜0.45 に下げる。
  \item \textbf{暗い所のノイズを少しでも減らしたい}: 0.65〜0.75 に上げる（継ぎ目が出ないか確かめる）。
\end{itemize}
\end{tip}

% ---------------------------------------------------------------------------------
\section{D. 飽和とみなす閾値}
センサーの値が\textbf{白飛び（飽和）の何割に達したら「白飛び」とみなして使わないか}を決めます。
白飛びの直前はセンサーの応答がまっすぐでないことがあるので、少し余裕をとっています。白飛びとみなした所とその隣は、必ず別のフレームから取ります。

\begin{tip}[変えるときの目安]
\begin{itemize}
  \item \textbf{ふだん}: 0.92（既定）。
  \item \textbf{明るい所の縁に色付きの縁取りや段差が出る}: 0.85〜0.88 に下げる（白飛びの手前をより早く避ける）。
  \item 0.95 以上に上げるのはおすすめしません（白飛び直前の、色の崩れた値を使うおそれがあります）。
\end{itemize}
\end{tip}

% ---------------------------------------------------------------------------------
\section{E. 重みのぼかし}
フレームを切り替える境目を、どのくらいの幅（画素）で\textbf{なめらかにならすか}を決めます。
ならすのは「明るいフレームを使う側」へだけで、白飛びした所に明るいフレームが混ざることはありません。

\begin{tip}[変えるときの目安]
\begin{itemize}
  \item \textbf{ふだん}: 16 px（既定）。
  \item \textbf{空や雲に、まだら模様・ちらつきが見える}: 32〜64 px に上げる。
  \item \textbf{大きなグラデーション全体をなめらかにしたい}: 128〜256 px まで上げられます。明るいフレームを使う範囲が少し広がります。
\end{itemize}
\end{tip}

% ---------------------------------------------------------------------------------
\section{F. 複数のフレームを平均してノイズを減らす}\label{sec:average}
オンにすると、それぞれの場所で\textbf{白飛びしていないフレームを、ノイズの少なさに応じた重みで平均}します（平均型の合成）。
オフにすると、場所ごとに「白飛びしていない中で最も明るい 1 枚」だけを使います（切り替え型）。
\begin{itemize}
  \item \textbf{既定}: オン。暗部・中間調のノイズが減ります。
  \item 効果は\textbf{露出の間隔が狭いほど大きく}なります。1⅔ 段刻みの 7 枚では、暗部のノイズ（分散）がおよそ 3 割減りました。2⅔ 段刻みでは 1 割ほどです。
  \item 減った割合は解析の結果に「平均型の合成: 暗部のノイズの分散 ○\%」と出ます。書き出す DNG のノイズの情報も、減った分だけ小さくします（Lightroom のノイズ軽減が効きすぎないように）。
  \item 動いた物は、次のゴースト対策（G）が平均から外すので、二重になりません。
  \item \textbf{オフにするとき}: 手持ちで撮って位置がわずかに合っていない写真で、細部がぼけて見えるとき。
\end{itemize}

% ---------------------------------------------------------------------------------
\section{G. 動いた物のゴーストを取り除く・動きの検出の感度}\label{sec:deghost}
撮影の間に動いた物（人・車・枝・波・雲・光の筋など）が、\textbf{二重や半透明に写る（ゴースト）}のを防ぎます。
フレームどうしを比べ、ノイズでは説明できない違いがある所を「動いた所」とし、そこでは\textbf{基準フレームと食い違うフレームを使いません}
（基準フレームが暗すぎる所では、それより明るいフレームを手本にします）。見つかった所は、表示の\ui{動いた所}で赤く示されます。
\begin{itemize}
  \item \textbf{既定}: オン、感度 50\%。
  \item 動いた所は、基準フレームを撮った瞬間の姿になります。波は長い露出の滑らかさが減って、波の形が見えるようになり、ノイズが少し増えることがあります。
  \item \textbf{感度}を上げると小さな違いも「動いた」とみなします。下げると、はっきり動いた所だけを扱います。
\end{itemize}

\begin{tip}[変えるときの目安]
\begin{itemize}
  \item \textbf{ふだん}: オン・50\%（既定）。
  \item \textbf{まだゴーストが残る}: 感度を 60〜75\% に上げる。
  \item \textbf{動いていない所まで赤く出る、または波・水面を長い露出のように滑らかにしたい}: 感度を 25\% に下げるか、オフにする。
\end{itemize}
\end{tip}

% ---------------------------------------------------------------------------------
\section{H. 明暗差の圧縮}\label{sec:compress}
このアプリのいちばん大事な設定です。\textbf{明るすぎる所（月・光源）を抑え、暗すぎる所（岩・草地）を持ち上げて、明暗差を縮めた DNG を作ります。}
暗室の「覆い焼き・焼き込み」にあたる処理です。輪郭に沿って行うので、月の縁や崖の輪郭ににじみ（ハロー）が出にくくなっています。
月の模様や岩の質感のような細かい模様は残ります。

\paragraph{なぜ必要？}
明暗差が 10 段を超えるような場面では、圧縮しない DNG を Lightroom で開くと、暗い所を持ち上げるために露光量を上げることになり、
そうすると月や光源が白く飛んで、ハイライトのスライダーを一番下げても模様が戻りません。
圧縮しておくと、\textbf{露光量を動かさずに、ハイライト・シャドウのスライダーだけで仕上げられる}ようになります。

\begin{figure}[htbp]
\centering
\begin{tikzpicture}
  \foreach \f/\t/\i in {r_col_0/{0\%（圧縮しない）}/0,r_col_50/{50\%（既定）}/1,r_col_100/{100\%（最大）}/2}{
    \node[inner sep=0,anchor=north west] (a\i) at ({\i*0.335*\linewidth},0) {\includegraphics[width=0.325\linewidth]{fig/\f.jpg}};
    \node[inner sep=0,anchor=north] (m\i) at ($(a\i.south)+(0,-0.15)$) {\includegraphics[width=0.16\linewidth]{fig/\f_moon.jpg}};
    \node[anchor=south,font=\sffamily\small\bfseries] at (a\i.north) {\t};
  }
\end{tikzpicture}
\caption{明暗差の圧縮の強さと、Lightroom（Camera Raw）で\textbf{何も調整せずに}開いたときの見え方。下は月の拡大。
0\% では全体が暗く、50\% では岩と空が自然な明るさで開き、100\% では岩がさらに明るく、月の模様（海）がはっきり出ます。}
\end{figure}

\begin{tabularx}{\linewidth}{l X}
\toprule
強さ & 内容 \\ \midrule
0\% & 圧縮しない。純粋な HDR（明るさの関係を一切変えない）。仕上げはすべて Lightroom で行う。明暗差の大きい場面では、露光量・ハイライトの調整が大変 \\
50\%（既定） & 標準の強さ。開いたときに暗部が見え、ハイライト・シャドウのスライダーで仕上げられる幅に収まる \\
100\% & 最大。明るい所をさらに下げ、暗い所をさらに持ち上げる。月の模様がはっきり出るが、いかにも「HDR」らしい平坦な見え方になりやすく、持ち上げた暗部のノイズも目立つ \\
\bottomrule
\end{tabularx}

\begin{tip}[変えるときの目安]
\begin{itemize}
  \item \textbf{ふだん}: 50\%（既定）。
  \item \textbf{月や光源の模様をもっと出したい／明暗差が極端（夜の灯台・月と地上）}: 75〜100\%。
  \item \textbf{自然な見え方を優先、または自分で一から仕上げたい}: 0〜30\%。
\end{itemize}
\end{tip}

\paragraph{強さを場面に合わせて自動で決める}
オンにすると、場面の明暗差（大まかな明るさの幅）から強さを決めます。合成した後、決めた強さがスライダーに表示されます。
\begin{center}
\begin{tabular}{lccc}
\toprule
大まかな明るさの幅 & 8 段以下 & 14 段 & 22 段以上 \\ \midrule
決まる強さ & 15\% & 50\% & 60\% \\
\bottomrule
\end{tabular}
\end{center}
既定はオフ（いつも 50\%）です。作例の 2 場面（柱状節理と月・夜の灯台）では、自動にすると 50\% と 59\% になります。

\paragraph{明るい所を抑える・暗い所を持ち上げる}
圧縮を、明るい側と暗い側で別々に強めたり弱めたりします（上の強さに対する割合、0〜200\%）。どちらも 100\% なら、上の強さのとおりです。
\begin{itemize}
  \item \textbf{持ち上げた暗部のノイズが気になる}: \ui{暗い所を持ち上げる}を 50\% ほどに下げる。
  \item \textbf{月や光源の模様をもっと出したいが、暗部は持ち上げたくない}: \ui{明るい所を抑える}を 150〜200\% に上げ、\ui{暗い所を持ち上げる}は 100\% のまま。
  \item 0\% にすると、その側は圧縮しません。
\end{itemize}

\begin{memo}[圧縮したときの DNG について]
\begin{itemize}
  \item 開いたときの明るさも自動で整えます（Lightroom で露光量を動かさなくても、ちょうどよい明るさで開きます）。
  \item 圧縮したときの DNG は、自動で LinearRaw 形式（K）になります。
  \item 強さは DNG の情報（XMP）に記録されます。
  \item Lightroom の\ui{かすみの除去}は、圧縮したときによく効きます（0\% のときは、月などがとても明るいため効きにくくなります）。
\end{itemize}
\end{memo}

% ---------------------------------------------------------------------------------
\section{I. 色の縁取り（倍率色収差）を補正}
レンズの性質で、明るい物の縁（特に画面の端の方）に\textbf{赤・緑・青の縁取り}が出ることがあります（倍率色収差）。
オンにすると、合成の直後に色ずれの量を写真から測って補正します。見積もりが確かな色だけを補正します。
\begin{itemize}
  \item \textbf{既定}: オン。普通はオンのままで構いません。
  \item 補正した量は「解析の結果」に「色収差の補正: 隅で R …px・B …px」と出ます。
  \item Lightroom の\ui{色収差を除去}は、明暗差を圧縮した DNG では十分に効かないため、このアプリで補正しておきます。
\end{itemize}

% ---------------------------------------------------------------------------------
\section{J. HDR 結合のボタン}
合成を始めるボタンです。状態によって表示が変わります。

\noindent\begin{tabularx}{\linewidth}{l X}
\toprule
表示 & 意味 \\ \midrule
\ui{HDR結合開始} & まだ合成していない。押すと合成する \\
\ui{結合済み}（灰色） & 今の設定で合成済み。そのまま書き出せる \\
\ui{再結合中…} & 合成した後に設定・フレーム・位置合わせを変えたので、自動で合成し直している（終わるまで書き出せない）。押すとすぐに始める \\
\bottomrule
\end{tabularx}

% ---------------------------------------------------------------------------------
\section{K. 書き出しの形式}
DNG の中身の形を選びます。\textbf{\ui{自動（おすすめ）}のままにしてください。}

\noindent\begin{tabularx}{\linewidth}{l X}
\toprule
形式 & 内容 \\ \midrule
自動（おすすめ） & 場面の明暗差から、下の 2 つのどちらかを自動で選ぶ。選んだ形式と理由が下に小さく出る。明暗差の圧縮を使うときは LinearRaw \\
CFA（色補間前） & 普通の RAW と同じ形。色補間（デモザイク）を Lightroom が行う。ただし Camera Raw は白から下およそ 16 段までしか扱えないので、明暗差がとても大きいと暗部に段差が出ることがある \\
LinearRaw（色補間済み） & このアプリが色補間した形。Camera Raw は小数のまま処理するので、明暗差の制限がない \\
\bottomrule
\end{tabularx}

% ---------------------------------------------------------------------------------
\section{L. レンズ補正を有効にして開く}
オンにすると、DNG を開いたときに Lightroom のレンズプロファイル補正（歪み・周辺光量）が最初から入った状態になります。
\begin{itemize}
  \item \textbf{既定}: オン。
  \item オンのときは、Lightroom の「既定の現像設定」（自分で決めたプロファイルなど）が使われず、Adobe の既定で開きます。自分の既定の設定で開きたいときはオフにしてください。
  \item オフでも、Lightroom の\ui{レンズ補正}→\ui{プロファイル補正を使用}で、レンズは自動で認識されます。
\end{itemize}

% =====================================================================================
\chapter{解析の結果の読み方}\label{ch:info}
% =====================================================================================

右下の\ui{解析の結果}には、合成についての情報が出ます。うまくいかないときの手がかりになります。

\begin{figure}[htbp]
\centering
\shot[0.55\linewidth]{c_info.jpg}{
  \node[num] at (-0.05,{1-0.088}) {1};
  \node[num] at (-0.05,{1-0.125}) {2};
  \node[num] at (-0.05,{1-0.30}) {3};
  \node[num] at (-0.05,{1-0.53}) {4};
  \node[num] at (-0.05,{1-0.615}) {5};
  \node[num] at (-0.05,{1-0.662}) {6};
  \node[num] at (-0.05,{1-0.712}) {7};
  \node[num] at (-0.05,{1-0.85}) {8};
}
\caption{解析の結果}
\end{figure}

\begin{enumerate}[label=\protect\cnum{\arabic*}]
  \item \textbf{カメラとレンズ}: 元の RAW から読み取ったもの。
  \item \textbf{Adobe DNG Converter}: 「あり」なら色と明るさを Camera Raw にそろえます。「なし」でも動きます。
  \item \textbf{露出比（隣どうし、全体で最適化した値）}: 隣り合うフレームの明るさの比。写真の中身から測り、2 枚離れた組の比も使って全体で矛盾しないように整えた値。カッコ内は EXIF の名目値（1/3 段刻みなど）とのずれ。「名目値」とあるのは、測れなかったので EXIF の値を使ったもの。
  \item \textbf{基準}: 基準フレームの番号と、一番暗い 1 枚より何段明るいか。その下に、平均型の合成で減った暗部のノイズ（F）と、動いた所の数と面積（G）が出ます。
  \item \textbf{色収差の補正}: 画面の隅での R・B のずれの補正量（画素）。
  \item \textbf{明暗差の圧縮}: 掛けた明るさの倍率の範囲（段）と、開いたときの明るさを整えた量。
  \item \textbf{最も暗いフレームでも飽和}: 一番暗い 1 枚でも白飛びしていた面積の割合。0\% でなければ、その部分は白飛びしたまま（もっと暗い 1 枚が必要）。
  \item \textbf{使った割合（面積）}: 各フレームが合成に使われた面積の割合。
\end{enumerate}

% =====================================================================================
\chapter{Lightroom での仕上げ方}\label{ch:lr}
% =====================================================================================

書き出した DNG は普通の RAW と同じように扱えます。明暗差の圧縮（既定 50\%）を使った場合の、おすすめの手順です。

\begin{enumerate}
  \item \textbf{レンズ補正}: \ui{レンズ補正}→\ui{プロファイル補正を使用}をオン（レンズは自動で認識されます）。
  \item \textbf{ホワイトバランス}: 撮影時の値で開きます。好みで調整します。
  \item \textbf{露光量}: 基本はそのまま（開いたときの明るさは整えてあります）。全体の明るさを少し変えたいときだけ動かします。
  \item \textbf{ハイライト・シャドウ}: ハイライトを下げると月や光源の模様が、シャドウを上げると暗部が見えてきます。
  \item \textbf{白レベル・黒レベル}: 明るい所・暗い所の締まりを整えます。
  \item \textbf{かすみの除去・テクスチャ・明瞭度}: 好みで。かすみの除去は少し（+10〜+20）でも効果があります。
  \item \textbf{ノイズ軽減}: 持ち上げた暗部のざらつきが気になるときは、\ui{ディテール}→\ui{ノイズ軽減}の\ui{輝度}を上げます（色のノイズは既定でほぼ消えます）。
\end{enumerate}

\begin{figure}[htbp]
\centering
\begin{tikzpicture}
  \node[inner sep=0] (a) {\includegraphics[width=0.48\linewidth]{fig/r_col_50.jpg}};
  \node[inner sep=0,right=0.02\linewidth of a] (b) {\includegraphics[width=0.48\linewidth]{fig/r_col_50_edit.jpg}};
  \node[above=2pt of a,font=\sffamily\small] {開いたまま（何も調整していない）};
  \node[above=2pt of b,font=\sffamily\small] {仕上げた後};
\end{tikzpicture}
\caption{明暗差の圧縮 50\% の DNG を Camera Raw で開いたもの（左）と、ハイライト $-60$・シャドウ $+60$・かすみの除去 $+10$・自然な彩度 $+15$ で仕上げたもの（右）。}
\end{figure}

\begin{figure}[htbp]
\centering
\begin{tikzpicture}
  \node[inner sep=0] (a) {\includegraphics[width=0.48\linewidth]{fig/r_lh_single.jpg}};
  \node[inner sep=0,right=0.02\linewidth of a] (b) {\includegraphics[width=0.48\linewidth]{fig/r_lh_50.jpg}};
  \node[above=2pt of a,font=\sffamily\small] {1 枚の RAW（真ん中の明るさ）};
  \node[above=2pt of b,font=\sffamily\small] {5 枚を合成（圧縮 50\%、何も調整していない）};
\end{tikzpicture}
\caption{夜の灯台の作例。合成すると、灯室の中と、草地・空の両方が写ります。}
\end{figure}

\begin{memo}[明暗差の圧縮を 0\% にしたとき]
圧縮しない DNG は、基準フレームと同じ明るさで開きます。暗い所を見せるには露光量を上げ、明るい所はハイライト・白レベルを大きく下げて調整してください。
明暗差がとても大きい場面では、この調整だけでは仕上げきれないことがあります。そのときは圧縮を使ってください。
\end{memo}

% =====================================================================================
\chapter{困ったときは}\label{ch:faq}
% =====================================================================================

\newcommand{\faq}[2]{\begin{tcolorbox}[colback=white,colframe=gray!60,title={\faQuestionCircle\ #1},fonttitle=\bfseries\sffamily,coltitle=black,colbacktitle=gray!12,boxrule=0.5pt,arc=2pt,left=6pt,right=6pt,breakable]#2\end{tcolorbox}}

\faq{アプリが開けない（「開発元を確認できない」と出る）}{
初めて開くときに macOS が止めています。ディスクイメージの\textsf{かんたんインストーラ.scpt}で入れ直すか、第 3 章の「手で入れる」の手順で開いてください。}

\faq{\ui{HDR結合開始}が押せない}{
\begin{itemize}
  \item RAW が 2 枚以上必要です。
  \item 別のカメラの RAW や、画像の大きさの違う RAW が混ざっていると合成できません（状態の表示に理由が出ます）。
  \item 処理中（読み込み・合成・書き出し）は押せません。終わるまで待ってください。
\end{itemize}}

\faq{\ui{DNG を書き出す…}が押せない}{
まだ合成していないか、設定を変えたので合成し直している途中です。最初は\ui{HDR結合開始}を押し、合成し直している途中なら\ui{結合済み}になるまで待ってください。}

\faq{つなぎ目に段差や、まだら模様が見える}{
\begin{itemize}
  \item \ui{重みのぼかし}を 32〜64 px に上げる。
  \item \ui{切り替えを始める明るさ}を 0.40〜0.45 に下げる。
  \item 三脚が動いた、手持ちで撮った場合は\ui{位置合わせ}を\ui{自動}にする。
\end{itemize}
\ui{前回との違い}の表示で、設定を変えた効果を確かめられます。}

\faq{明るい所の縁に色の縁取りが出る}{
\begin{itemize}
  \item \ui{色の縁取り（倍率色収差）を補正}がオンか確かめる。
  \item \ui{飽和とみなす閾値}を 0.85〜0.88 に下げる。
  \item 撮影中に動いたもの（雲・人・枝）の縁なら、\ui{動いた物のゴーストを取り除く}がオンか確かめ、感度を上げてみる（\ref{sec:deghost}節）。
\end{itemize}}

\faq{暗い所のノイズが目立つ}{
\begin{itemize}
  \item \ui{複数のフレームを平均してノイズを減らす}がオンか確かめる。
  \item 明暗差の圧縮を下げるか、\ui{暗い所を持ち上げる}を下げる（暗部を持ち上げる量が減ります）。
  \item Lightroom の\ui{ノイズ軽減}→\ui{輝度}を上げる。
  \item 撮影時に、もっと明るい 1 枚（長い露出）を加える。
\end{itemize}}

\faq{動いたもの（人・車・波）が二重に写る}{
\begin{itemize}
  \item \ui{動いた物のゴーストを取り除く}をオンにする（既定でオン）。表示の\ui{動いた所}で、見つかった所が赤く出ます。
  \item 赤く出ていない所に二重写りが残るときは、\ui{動きの検出の感度}を 60〜75\% に上げる。
  \item それでも残るときは、動くものが主役の場面では枚数を減らして手早く撮ってください。
\end{itemize}}

\faq{波や水面が、長い露出の滑らかさにならない・ノイズが増えた}{
ゴースト対策が、動く水面を基準フレームの瞬間の姿にそろえています。滑らかにしたいときは、\ui{動きの検出の感度}を下げるか、\ui{動いた物のゴーストを取り除く}をオフにしてください。}

\faq{右の欄の下の方（書き出し・解析の結果）が見えない}{
ウインドウが低いと右の欄に収まりません。右の欄の上でスクロールしてください（ウインドウを大きくしても見えます）。}

\faq{Lightroom の「かすみの除去」が効かない}{
明暗差の圧縮が 0\% のときは、月や光源がとても明るいため、かすみの除去がほとんど効きません。圧縮を使ってください（既定の 50\% ならよく効きます）。}

\faq{プレビューと、DNG を Lightroom・Camera Raw で開いたときの見え方が少し違う}{
\begin{itemize}
  \item プレビューは、DNG を何も調整せずに開いたときに近づけていますが、Lightroom で選んだプロファイル（カメラのポートレートなど）や、自分の既定の現像設定による色・コントラストの違いは再現しません。
  \item Adobe DNG Converter が無いと、明るさの基準が Adobe と少し違います（入れるとそろいます）。
  \item \ui{ハイライトを見やすく}をオンにして明暗差の圧縮を 0\% にしているときは、表示だけの見やすさのため、開いたときとは違う見え方になります。
\end{itemize}}

\faq{色や明るさが、元の RAW を Lightroom で開いたときと少し違う}{
Adobe DNG Converter を入れると、そろいます（メニュー\ui{RawHDR Composer}→\ui{Adobe DNG Converter…}）。}

\faq{書き出した DNG のファイルが大きい}{
小数（浮動小数点）で値を持つため、1 枚あたり数十〜百数十 MB になります。}

% =====================================================================================
\appendix
\chapter{付録}
% =====================================================================================

\section{キーボードショートカット}
\begin{center}
\begin{tabular}{ll}
\toprule
キー & 操作 \\ \midrule
\key{⌘}\key{O} & RAW を開く（追加する） \\
\key{⌘}\key{R} & HDR 結合（合成する・合成し直す） \\
\key{⌘}\key{E} & DNG を書き出す \\
\key{⌘}\key{0} & プレビューの全体を表示 \\
\key{⌘}\key{1} & プレビューを 100\% で表示 \\
\key{⌘}\key{＋}／\key{⌘}\key{−} & 拡大／縮小 \\
\key{Return} & 青いボタン（\ui{HDR結合開始}・\ui{DNG を書き出す…}）を押す \\
\key{⌘}\key{W}／\key{⌘}\key{Q} & ウインドウを閉じる／アプリを終了 \\
\bottomrule
\end{tabular}
\end{center}

\section{用語}
\begin{description}[style=nextline,leftmargin=2em]
  \item[RAW] カメラのセンサーの値をほぼそのまま記録したファイル。後からホワイトバランスや明るさを大きく変えられる。
  \item[DNG] Adobe が公開している RAW の共通形式。このアプリは合成結果を DNG で書き出す。
  \item[HDR] ハイダイナミックレンジ。明るい所から暗い所まで、広い明るさの幅を持った写真。
  \item[AEB（オートブラケット）] シャッターを切るたびに、カメラが自動で明るさを変えて何枚か撮る機能。
  \item[段（EV）] 明るさの単位。1 段明るい＝光の量が 2 倍。
  \item[白飛び（飽和）] センサーが受けられる光の上限を超え、明るさの違いが記録されていない状態。
  \item[基準フレーム] DNG の明るさと撮影情報の基準にする 1 枚。
  \item[由来マップ] 合成結果の各場所が、どのフレームから来たかを色で示した図。
  \item[CFA・LinearRaw] DNG の中身の形。CFA は色補間（デモザイク）前、LinearRaw は色補間済み。
  \item[倍率色収差] レンズの性質で、色ごとに像の大きさがわずかに違い、明るい物の縁に色の縁取りが出ること。
  \item[ゴースト] 撮影の間に動いた物が、合成で二重や半透明に写ること。
  \item[平均型の合成] 白飛びしていない複数のフレームを、ノイズの少なさに応じた重みで平均すること。ノイズが減る。
\end{description}

\section{設定の早見表}
\begin{center}
\begin{tabularx}{\linewidth}{l c X}
\toprule
設定 & 既定値 & こんなときに変える \\ \midrule
切り替えを始める明るさ & 0.55 & 継ぎ目・段差 → 0.40〜0.45 ／ 暗部のノイズを減らしたい → 0.65〜0.75 \\
飽和とみなす閾値 & 0.92 & 明るい所の縁に色・段差 → 0.85〜0.88 \\
重みのぼかし & 16 px & まだら・ちらつき → 32〜64 px \\
平均してノイズを減らす & オン & 手持ちで細部がぼける → オフ \\
ゴーストを取り除く・感度 & オン・50\% & 二重写りが残る → 60〜75\% ／ 水面を滑らかに → 25\% かオフ \\
明暗差の圧縮 & 50\% & 月・光源の模様を出す → 75〜100\% ／ 自然に・自分で仕上げる → 0〜30\% ／ 迷ったら「自動」 \\
明るい所を抑える・暗い所を持ち上げる & 100\% & 暗部のノイズが気になる → 暗い所を 50\% \\
色の縁取りを補正 & オン & 普通はオンのまま \\
書き出しの形式 & 自動 & 普通は自動のまま \\
\bottomrule
\end{tabularx}
\end{center}

\end{document}
